\documentclass[12pt,a4paper]{ctexart}

\usepackage[margin=2.5cm]{geometry}
\usepackage{listings}
\usepackage{xcolor}
\usepackage{hyperref}
\usepackage{booktabs}
\usepackage{enumitem}

\hypersetup{
  colorlinks=true,
  linkcolor=blue,
  urlcolor=blue
}

% 代码块样式
\definecolor{codebg}{RGB}{245,245,245}
\definecolor{codeframe}{RGB}{200,200,200}
\definecolor{codecomment}{RGB}{100,100,100}
\definecolor{codekeyword}{RGB}{0,0,160}

\lstdefinestyle{shell}{
  backgroundcolor=\color{codebg},
  frame=single,
  rulecolor=\color{codeframe},
  basicstyle=\ttfamily\small,
  breaklines=true,
  columns=fullflexible,
  keepspaces=true,
  showstringspaces=false,
  commentstyle=\color{codecomment},
  keywordstyle=\color{codekeyword},
  language=bash
}

\lstdefinestyle{ini}{
  backgroundcolor=\color{codebg},
  frame=single,
  rulecolor=\color{codeframe},
  basicstyle=\ttfamily\small,
  breaklines=true,
  columns=fullflexible,
  keepspaces=true,
  showstringspaces=false,
  commentstyle=\color{codecomment}
}

\lstdefinestyle{log}{
  backgroundcolor=\color{codebg},
  frame=single,
  rulecolor=\color{codeframe},
  basicstyle=\ttfamily\footnotesize,
  breaklines=true,
  columns=fullflexible,
  keepspaces=true,
  showstringspaces=false
}

\title{RuoYi CMDB 项目 systemd 开机自启配置文档}
\author{运维配置记录}
\date{\today}

\begin{document}

\maketitle

\tableofcontents
\newpage

\section{环境概述}

\begin{itemize}[leftmargin=2em]
  \item 操作系统：Ubuntu 26.04 LTS（systemd 版本 259.x）
  \item 应用框架：RuoYi（基于 Spring Boot 2.2.13.RELEASE）
  \item 应用 jar：\texttt{/opt/cmdb/admin/cmdb-admin-1.0.0.jar}
  \item 工作目录：\texttt{/opt/cmdb/admin}
  \item 外部配置：\texttt{/opt/cmdb/admin/config/application-prod.yml}
  \item 日志目录：\texttt{/opt/cmdb/admin/logs}
  \item 依赖服务：\texttt{redis-server.service}（Ubuntu 上 Redis 的服务名）
  \item 应用端口：8089，context path 为 \texttt{/cmdb}
\end{itemize}

\section{创建 systemd 服务文件}

在 \texttt{/etc/systemd/system/} 下创建服务文件：

\begin{lstlisting}[style=shell]
sudo vim /etc/systemd/system/cmdb.service
\end{lstlisting}

写入以下内容：

\begin{lstlisting}[style=ini]
[Unit]
Description=RuoYi CMDB Backend Service
After=network.target redis-server.service
Wants=redis-server.service

[Service]
Type=simple
User=root
WorkingDirectory=/opt/cmdb/admin
ExecStart=/usr/bin/java -jar /opt/cmdb/admin/cmdb-admin-1.0.0.jar --spring.profiles.active=prod --spring.config.additional-location=file:/opt/cmdb/admin/config/application-prod.yml --logging.file.path=/opt/cmdb/admin/logs
SuccessExitStatus=143
Restart=on-failure
RestartSec=10s

[Install]
WantedBy=multi-user.target
\end{lstlisting}

\subsection{关键配置说明}

\begin{itemize}[leftmargin=2em]
  \item \texttt{After=network.target redis-server.service}：确保 cmdb 在网络和 Redis 之后启动。
  \item \texttt{Wants=redis-server.service}：cmdb 启动时主动触发 Redis 启动；用 \texttt{Wants} 而非 \texttt{Requires}，避免 Redis 故障时强依赖导致 cmdb 直接失败。
  \item \texttt{WorkingDirectory}：指定工作目录，避免相对路径问题。
  \item \texttt{ExecStart}：使用 java 绝对路径，显式指定 prod profile、外部配置文件与日志路径。
  \item \texttt{SuccessExitStatus=143}：Java 应用被正常停止时退出码为 143，视为正常退出。
  \item \texttt{Restart=on-failure} + \texttt{RestartSec=10s}：异常退出后自动重试。
  \item \textbf{不显式设置} \texttt{StandardOutput}/\texttt{StandardError}：控制台输出默认进入 journal，用 \texttt{journalctl} 查看。
\end{itemize}

\section{日志策略}

系统日志分为两类，分别管理：

\begin{table}[h]
\centering
\begin{tabular}{@{}p{3.2cm}p{4.2cm}p{2.6cm}p{3.5cm}@{}}
\toprule
日志类型 & 位置 & 自动清理 & 查看方式 \\
\midrule
RuoYi 业务日志（Logback） & \texttt{/opt/cmdb/admin/logs} & 取决于 \texttt{maxHistory} 配置 & \texttt{tail -f} 或直接查看文件 \\
systemd 控制台日志 & journald 内部 & 是（默认轮转） & \texttt{journalctl -u cmdb} \\
\bottomrule
\end{tabular}
\end{table}

\paragraph{RuoYi 业务日志。}
RuoYi 通过 \texttt{logback-spring.xml} 配置日志滚动。需确认其中是否设置了 \texttt{<maxHistory>}（例如 30 天），否则旧日志会持续堆积，不会自动删除。

\paragraph{systemd 控制台日志。}
journald 默认按磁盘占用和文件大小自动轮转、清理，无需额外配置。可用 \texttt{journalctl --disk-usage} 查看占用，或用 \texttt{journalctl --vacuum-time=7d} 手动清理。

\section{启用与启动}

\begin{lstlisting}[style=shell]
# 创建日志目录
sudo mkdir -p /opt/cmdb/admin/logs

# 重新加载 systemd 配置
sudo systemctl daemon-reload

# 设置开机自启
sudo systemctl enable cmdb

# 启动服务
sudo systemctl start cmdb

# 查看状态
sudo systemctl status cmdb
\end{lstlisting}

\section{启动验证结果}

启动日志显示应用已成功初始化，关键信息如下：

\begin{lstlisting}[style=log]
The following profiles are active: prod
dynamic-datasource - load a datasource named [db1] success
Tomcat initialized with port(s): 8089 (http)
Tomcat started on port(s): 8089 (http) with context path '/cmdb'
Started AdminApplication in 18.667 seconds (JVM running ...)
\end{lstlisting}

由此可确认：

\begin{itemize}[leftmargin=2em]
  \item \texttt{prod} profile 生效；
  \item 动态数据源 \texttt{db1} 加载成功；
  \item Tomcat 在 8089 端口启动，context path 为 \texttt{/cmdb}；
  \item 应用在 18.667 秒内完成启动；
  \item Redis 作为前置依赖已先行启动。
\end{itemize}

\subsection{端到端验证命令}

\begin{lstlisting}[style=shell]
# 服务状态
systemctl status cmdb

# 业务日志目录
ls -l /opt/cmdb/admin/logs/

# 访问接口（返回 200/302/401 均表示链路通畅）
curl -I http://127.0.0.1:8089/cmdb/

# 确认启动完成标志
journalctl -u cmdb | grep "Started AdminApplication"
\end{lstlisting}

\section{已知警告与说明}

启动日志中出现以下警告，不影响当前运行：

\begin{itemize}[leftmargin=2em]
  \item \texttt{WARNING: A restricted method in java.lang.System has been called}（Tomcat native 调用受限方法）；
  \item \texttt{WARNING: sun.misc.Unsafe::objectFieldOffset}（Netty 调用已弃用 API）。
\end{itemize}

这些警告源于较新的 JDK 对旧版 Spring Boot 2.2.13 / Tomcat 9.0.41 依赖的兼容性提示。当前可忽略；若未来升级到更高 JDK 版本（如 JDK 25+），受限方法可能被彻底禁用，届时需同步升级 Spring Boot 与相关依赖。

\section{常用运维命令}

\begin{lstlisting}[style=shell]
# 启动 / 停止 / 重启
sudo systemctl start cmdb
sudo systemctl stop cmdb
sudo systemctl restart cmdb

# 开机自启开关
sudo systemctl enable cmdb
sudo systemctl disable cmdb

# 查看实时控制台日志
journalctl -u cmdb -f

# 查看最近 100 行
journalctl -u cmdb -n 100

# 确认 Redis 服务名与状态
systemctl list-units --type=service | grep -i redis
\end{lstlisting}

\end{document}